\section{Bab II: Sejarah Penyusunan dan Pengesahan UUD NRI Tahun 1945}
\subsection{A. Sejarah Penyusunan oleh BPUPK}
Undang-Undang Dasar Negara Indonesia merdeka dirancang sejak tanggal 29 Mei 1945 sampai tanggal 16 Juli 1945 oleh Badan Penyelidik Usaha-Usaha Persiapan Kemerdekaan (BPUPK). Terdapat tiga masalah pokok dalam rancangan UUD yang disampaikan oleh Sukarno, yaitu: Pernyataan tentang Indonesia merdeka, Pembukaan UUD, dan Batang Tubuh UUD. Dalam perumusannya, dibentuk tiga panitia:
\begin{itemize}
	\item Panitia Perancang Hukum Dasar yang diketuai oleh Ir. Sukarno.
	\item Panitia Perancang Keuangan dan Ekonomi yang diketuai oleh Drs. Mohammad Hatta.
	\item Panitia Perancang Pembelaan Tanah Air yang diketuai oleh Abikoesno Tjokrosoejoso.
\end{itemize}
Panitia Perancang UUD kemudian membentuk Panitia Kecil yang diketuai oleh Prof. Dr. Mr. R. Supomo pada 11 Juli 1945, serta Panitia Penghalus Bahasa (Husein Djajadiningrat, Agoes Salim, dan Supomo).

\subsection{B. Pengesahan oleh PPKI dan Peran Tokoh Pemersatu}
BPUPK dibubarkan pada 7 Agustus 1945 dan digantikan oleh Panitia Persiapan Kemerdekaan Indonesia (PPKI). Pada sidang pertama tanggal 18 Agustus 1945 (diperingati sebagai Hari Konstitusi), PPKI menetapkan putusan penting: mengesahkan UUD NRI Tahun 1945, memilih Ir. Sukarno sebagai Presiden dan Drs. Mohammad Hatta sebagai Wakil Presiden, serta membentuk Komite Nasional.

\textbf{Peran Mr. Kasman Singodimedjo:}
Dalam proses pengesahan ini, Mr. Kasman Singodimedjo memiliki peran penting sebagai tokoh pemersatu bangsa. Beliau dipercaya oleh Sukarno dan Mohammad Hatta untuk meluluhkan hati perwakilan umat Islam (seperti Ki Bagoes Hadikoesoemo) agar menyetujui penghapusan tujuh kata dalam Piagam Jakarta menjadi \textit{"Ketuhanan Yang Maha Esa"}. Hal ini dilakukan semata-mata demi menjaga keutuhan Negara Kesatuan Republik Indonesia (NKRI). Atas pengabdiannya, beliau diangkat sebagai Pahlawan Nasional pada tahun 2018.